% !TeX root = ../my-thesis.tex

% 中英文摘要和关键字

\begin{abstract}
  近年来，互联网行业从高速增长期迈入深度调整期。在经济下行、增量市场收窄与大模型技术冲击的多重压力下，沿用"烟囱式"职能分工的传统研发组织架构暴露出高协作成本、低响应速度、技术负债持续积累等结构性缺陷，已成为制约互联网企业创新效能的核心矛盾。如何推动研发组织从职能型向产品型的根本性转型，成为互联网企业在新阶段维持技术竞争力的关键课题。

  本研究以互联网企业研发管理为研究对象，以百度公司为典型案例，围绕"组织重构—效能度量—技术资产治理"三条主线展开系统分析。理论框架层面，本研究在技术资产（TA）与技术负债（TD）分析框架的基础上，引入Team Topologies产品型组织设计理论、Conway's Law（康威定律）与DORA研发效能指标体系，构建了涵盖"组织形态—技术架构—研发效能"全链条的分析工具集。

  在百度案例分析中，通过PEST分析、波特五力模型、SWOT分析和AMO框架系统识别了研发管理的现状与问题；并以Team Topologies框架对百度TPG调整为BMU（基础模型与理解）与AMU（应用与多模态）双BU架构的组织演变进行结构性解读：BMU承担平台团队角色，专注技术资产的集中生产与标准化输出；AMU承担流对齐团队集合角色，以业务价值流为核心实现快速交付；技术委员会承担赋能团队角色，负责跨BU技术标准传递。这一组织拓扑的演进，是对Conway's Law"反向定律"策略的主动实践，从根本上重塑了技术资产的生产与流动机制。本研究进一步引入DORA四项效能指标对转型成效进行量化分析，并提出技术资产流动效率（$\eta_{\mathrm{TA}}$）作为衡量组织转型效益的补充指标。

  研究结论表明：组织重构是效能提升的结构性前提，而非可选项；以DORA指标为基准建立常态化效能度量体系，可以为组织转型提供客观评估依据；将TA管理上升为战略级议题、建立TA全生命周期管理机制，是提升研发投入产出比的根本路径；激励机制改革需与产品型组织新形态相匹配，才能避免激励结构与组织结构相互抵消。本研究将工程实践领域的产品型组织理论引入管理学研究框架，并整合TA/TD治理视角，为中国互联网企业研发管理转型提供了理论指导与实践参考。

  % 关键词用"英文逗号"分隔，输出时会自动处理为正确的分隔符
  \thusetup{
    keywords = {互联网企业, 研发组织重构, 产品型组织, Team Topologies, 技术资产, 技术负债, 研发效能},
  }
\end{abstract}

\begin{abstract*}
  In recent years, the Internet industry has transitioned from a phase of rapid growth to a period of deep structural adjustment. Under the compounding pressures of economic headwinds, shrinking incremental markets, and the disruptive impact of large language model technologies, traditional R\&D organizations structured around functional silos have exposed critical structural deficiencies—high coordination costs, slow responsiveness, and the persistent accumulation of technical debt—that have become the core bottleneck constraining innovation efficiency. How to fundamentally transform R\&D organizations from functional to product-oriented structures has emerged as a pivotal challenge for Internet enterprises seeking to maintain technological competitiveness.

  This study takes R\&D management in Internet enterprises as its research subject and uses Baidu Inc. as a representative case, organized around three main threads: organizational restructuring, efficiency measurement, and technical asset governance. On the theoretical side, the study builds upon the Technical Assets (TA) and Technical Debt (TD) analytical framework and integrates the Team Topologies product-oriented organization design theory, Conway's Law, and the DORA (DevOps Research and Assessment) metrics system, constructing a comprehensive analytical toolkit spanning the full chain from organizational structure to technical architecture to R\&D efficiency.

  In the Baidu case analysis, PEST analysis, Porter's Five Forces, SWOT analysis, and the AMO framework are applied to systematically identify the current state and challenges of R\&D management. The study then employs the Team Topologies framework to provide a structural interpretation of Baidu's organizational evolution from TPG to the dual-BU architecture of BMU (Foundation Model and Understanding) and AMU (Applications and Multimodality): BMU assumes the role of a Platform Team, focusing on the centralized production and standardized delivery of technical assets; AMU functions as a collection of Stream-aligned Teams, delivering value directly through business-oriented workflows; the Technology Committee plays the role of an Enabling Team, responsible for cross-BU standards and best-practice dissemination. This topological evolution represents an active application of the ``Inverse Conway Maneuver,'' fundamentally reshaping the production and flow mechanisms of technical assets. The study further introduces DORA's four key metrics to quantify transformation outcomes, and proposes Technical Asset Flow Efficiency ($\eta_{\mathrm{TA}}$) as a supplementary indicator for evaluating the returns on organizational transformation.

  The research conclusions indicate that: organizational restructuring is a structural prerequisite for efficiency improvement, not an optional measure; establishing a normalized efficiency measurement system anchored by DORA metrics provides an objective basis for evaluating organizational change; elevating TA management to a strategic-level agenda and building a full lifecycle TA governance mechanism is the fundamental path to improving R\&D return on investment; and incentive mechanism reforms must align with the new product-oriented organizational structure to prevent incentive systems from counteracting organizational design. By introducing product-oriented organization theory from engineering practice into a management research framework, and integrating the TA/TD governance perspective, this study provides theoretical guidance and practical reference for the R\&D management transformation of Chinese Internet enterprises.

  % Use comma as separator when inputting
  \thusetup{
    keywords* = {Internet enterprise, R\&D organization restructuring, product-oriented organization, Team Topologies, technical assets, technical debt, R\&D efficiency},
  }
\end{abstract*}
